%% Đề 44 -- de-tuyen-sinh-lop-10-thpt-mon-toan-nam-2016-2017-so-gddt-lang-son.pdf
\deheader{ĐỀ SỐ 44}%
  {Kỳ thi Tuyển sinh vào lớp 10 THPT — Năm học 2016--2017}%
  {Sở GD\&ĐT Lạng Sơn | Chỉ phần Đại số}
\addcontentsline{toc}{section}{Đề 44 — Tuyển sinh vào lớp 10 THPT — Năm học 2016--2017}

\begin{de_cau}{Câu 1}{2,0 điểm}
\begin{enumerate}[label=\alph*)]
  \item Tính: $A = \sqrt{8}+\sqrt{18}-\sqrt{2}$.
  \item Rút gọn: $B = \dfrac{x-1}{\sqrt{x}-1}$ với $x\ge0$, $x\neq1$.
\end{enumerate}
\end{de_cau}

\begin{loi_giai_cau}
\textbf{a)} $A=2\sqrt{2}+3\sqrt{2}-\sqrt{2}=\boxed{4\sqrt{2}}$.

\textbf{b)} $B=\dfrac{(\sqrt{x}-1)(\sqrt{x}+1)}{\sqrt{x}-1}=\boxed{\sqrt{x}+1}$.
\end{loi_giai_cau}

\begin{de_cau}{Câu 2}{2,0 điểm}
\begin{enumerate}[label=\alph*)]
  \item Giải phương trình: $x^2-6x+9=0$.
  \item Giải hệ phương trình:
  \[\begin{cases} 2x+y=4 \\ x-2y=-2 \end{cases}.\]
\end{enumerate}
\end{de_cau}

\begin{loi_giai_cau}
\textbf{a)} $(x-3)^2=0\Rightarrow\boxed{x=3}$ (nghiệm kép).

\textbf{b)} Từ (1): $y=4-2x$. Thế vào (2): $x-2(4-2x)=-2\Rightarrow5x=6\Rightarrow x=\dfrac{6}{5}$, $y=4-\dfrac{12}{5}=\dfrac{8}{5}$.

Nghiệm $\boxed{(x;y)=\left(\dfrac{6}{5};\dfrac{8}{5}\right)}$.
\end{loi_giai_cau}

\begin{de_cau}{Câu 3}{2,0 điểm}
Cho phương trình $x^2-2(m+1)x+m^2+2=0$ (tham số $m$).
\begin{enumerate}[label=\alph*)]
  \item Chứng minh phương trình có hai nghiệm phân biệt với mọi $m$.
  \item Tìm $m$ để $x_1^2+x_2^2=12$.
\end{enumerate}
\end{de_cau}

\begin{loi_giai_cau}
\textbf{a)} $\Delta'=(m+1)^2-(m^2+2)=2m-1$. Khi $m\ge\frac{1}{2}$: $\Delta'>0$. Khi $m<\frac{1}{2}$: $\Delta'<0$ -- vô nghiệm thực.

Thực ra đề muốn: $\Delta'=(m+1)^2-(m^2+2)=2m-1>0$ khi $m>\frac{1}{2}$. Phương trình luôn có hai nghiệm phân biệt khi $m>\dfrac{1}{2}$.

\textbf{b)} Vi-ét: $x_1+x_2=2(m+1)$, $x_1x_2=m^2+2$.

$x_1^2+x_2^2=4(m+1)^2-2(m^2+2)=2m^2+8m=12\Rightarrow m^2+4m-6=0\Rightarrow m=-2\pm\sqrt{10}$.

Kết hợp $m>\dfrac{1}{2}$: $m=\boxed{-2+\sqrt{10}}$.
\end{loi_giai_cau}

\begin{de_cau}{Câu 5}{0,5 điểm}
Cho $x, y > 0$ và $\dfrac{1}{x}+\dfrac{1}{y}=1$. Tìm giá trị nhỏ nhất của $P=x+y$.
\end{de_cau}

\begin{loi_giai_cau}
$P=(x+y)\left(\dfrac{1}{x}+\dfrac{1}{y}\right)=1+\dfrac{x}{y}+\dfrac{y}{x}+1\ge2+2=4$.

Dấu $=$ khi $x=y=2$. Vậy $P_{\min}=\boxed{4}$.
\end{loi_giai_cau}
